\chapter{Տվյալների ֆորմատ և մոդել}

FPGA/PS տվյալները սերվեր ուղարկելուց առաջ ձևափոխվում են կառուցվածքային ֆորմատի։
Օրինակ՝

\begin{verbatim}
{
  "device": "AUP_ZU3_01",
  "experiment_id": 12,
  "sensor": "vibration",
  "timestamp": 1789063200,
  "sequence": 1524,
  "raw_value": 728,
  "processed_value": 681,
  "status": "normal"
}
\end{verbatim}

PS-ը այս տվյալները կարող է ուղարկել API-ին, օրինակ՝ \texttt{POST /api/v1/data}։

Տվյալների բազան պետք է աջակցի \textbf{փորձերի (experiments)} հասկացությունը, ոչ թե
պահի միայն առանձին արժեքներ։ Հայեցակարգային մոդելը.

\begin{itemize}[leftmargin=1.4em]
  \item \textbf{Experiment} — \texttt{experiment\_id}, \texttt{device\_id},
        \texttt{configuration}, \texttt{start\_time}, \texttt{end\_time}։
  \item \textbf{Sensor} — \texttt{sensor\_id}, \texttt{sensor\_type},
        \texttt{sample\_rate}, \texttt{unit}։
  \item \textbf{Sample} — \texttt{experiment\_id}, \texttt{sensor\_id},
        \texttt{timestamp}, \texttt{sequence\_number}, \texttt{raw\_value},
        \texttt{processed\_value}, \texttt{status}։
\end{itemize}

Այս մոդելը թույլ է տալիս հետագայում համեմատել փորձերը (օր.՝ Experiment~\#12 vs \#18)։
